%!TEX encoding = UTF8
%!TEX root = 

%%%%%%%%%%%%%%%%%%%%%%%%
%%%%%%%%%% SET 1 %%%%%%%%%%
%%%%%%%%%%%%%%%%%%%%%%%% 


\exe{}{
	Écrire les sommes et différences de rationnels sous forme de fraction irréductible.
	\begin{multicols}{3}
	\begin{enumerate}[itemsep=10pt]
		\item $1 - \dfrac12$
		\item $\dfrac78 + \dfrac{9}{40}$
		\item $\dfrac{3}{4} + \dfrac{2}{5}$
		\item $\dfrac{7}{6} - \dfrac{5}{8}$
		\item $\dfrac{9}{10} + \dfrac{2}{3} - \dfrac{1}{4}$
		\item $\dfrac{-5}{7} + \dfrac{8}{9}$
%		\item $\dfrac{4}{5} - \dfrac{3}{10} + \dfrac{-5}{6}$
%		\item $\dfrac{-1}{2} + \dfrac{7}{8}$
%		\item $\dfrac{11}{12} - \dfrac{5}{9} + \dfrac{7}{10}$
%		\item $\dfrac{2}{3} - \dfrac{5}{4}$
%		\item $\dfrac{-1}{3} + \dfrac{7}{9} - \dfrac{5}{6}$
%		\item $\dfrac{3}{8} - \dfrac{1}{2}$
	\end{enumerate}
	\end{multicols}
}{exe:rationnals}
{
	\begin{multicols}{2}
	\begin{enumerate}
		\item $1 - \dfrac{1}{2} = \dfrac{1}{2}$
		\item $\dfrac{7}{8} + \dfrac{9}{40} = \dfrac{11}{10}$
		\item $\dfrac{3}{4} + \dfrac{2}{5} = \dfrac{23}{20}$
		\item $\dfrac{7}{6} - \dfrac{5}{8} = \dfrac{13}{24}$
		\item $\dfrac{9}{10} + \dfrac{2}{3} - \dfrac{1}{4} = \dfrac{79}{60}$
		\item $\dfrac{-5}{7} + \dfrac{8}{9} = \dfrac{11}{63}$
%		\item $\dfrac{4}{5} - \dfrac{3}{10} + \dfrac{-5}{6} = \dfrac{-1}{3}$
%		\item $\dfrac{-1}{2} + \dfrac{7}{8} = \dfrac{3}{8}$
%		\item $\dfrac{11}{12} - \dfrac{5}{9} + \dfrac{7}{10} = \dfrac{191}{180}$
%		\item $\dfrac{2}{3} - \dfrac{5}{4} = \dfrac{-7}{12}$
%		\item $\dfrac{-1}{3} + \dfrac{7}{9} - \dfrac{5}{6} = \dfrac{-7}{18}$
%		\item $\dfrac{3}{8} - \dfrac{1}{2} = \dfrac{-1}{8}$
	\end{enumerate}
	\end{multicols}
}


\exe{}{
	Écrire les produits et divisions de rationnels sous forme de fraction irréductible.
	
	\begin{multicols}{3}
	\begin{enumerate}[itemsep=10pt]
		\item $\dfrac{\ \ -3 \ \ }{\dfrac54}$
		\item $\dfrac{63}{20} \times \dfrac{15}{14}$
		\item $\dfrac{3}{4} \times \dfrac{2}{5}$
		\item $\dfrac{\ \ \dfrac{7}{6}\ \ }{\dfrac{5}{8}}$
		\item $\dfrac{9}{10} \times \dfrac{2}{3} \times \dfrac{4}{1}$
		\item $\dfrac{-5}{7} \times \dfrac{8}{9}$
%		\item $\left(\dfrac{4}{5}\right)^{-1} \times \dfrac{3}{10}$
%		\item $\dfrac{-1}{2} \times \left(\dfrac{7}{8}\right)^{-1}$
%		\item $\dfrac{11}{12} \times \dfrac{5}{9} \times \left(\dfrac{7}{10}\right)^{-1}$
%		\item $\dfrac{2}{3} \times \dfrac{5}{4}$
%		\item $\dfrac{-1}{3} \times \left(\dfrac{7}{9}\right)^{-1}$
%		\item $\dfrac{3}{8} \times \dfrac{1}{2}$
	\end{enumerate}
	\end{multicols}
}{exe:rationnals2}
{
	\begin{multicols}{2}
	\begin{enumerate}[itemsep=10pt]
		\item $\dfrac{-3}{\dfrac{5}{4}} = \dfrac{-12}{5}$
		\item $\dfrac{63}{20} \times \dfrac{15}{14} = \dfrac{27}{8}$
		\item $\dfrac{3}{4} \times \dfrac{2}{5} = \dfrac{3}{10}$
		\item $\dfrac{\ \ \dfrac{7}{6}\ \ }{\dfrac{5}{8}} = \dfrac{28}{15}$
		\item $\dfrac{9}{10} \times \dfrac{2}{3} \times \dfrac{4}{1} = \dfrac{12}{5}$
		\item $\dfrac{-5}{7} \times \dfrac{8}{9} = \dfrac{-40}{63}$
%		\item $\left(\dfrac{4}{5}\right)^{-1} \times \dfrac{3}{10} = \dfrac{3}{8}$
%		\item $\dfrac{-1}{2} \times \left(\dfrac{7}{8}\right)^{-1} = \dfrac{-4}{7}$
%		\item $\dfrac{11}{12} \times \dfrac{5}{9} \times \left(\dfrac{7}{10}\right)^{-1} = \dfrac{275}{378}$
%		\item $\dfrac{2}{3} \times \dfrac{5}{4} = \dfrac{5}{6}$
%		\item $\dfrac{-1}{3} \times \left(\dfrac{7}{9}\right)^{-1} = \dfrac{-3}{7}$
%		\item $\dfrac{3}{8} \times \dfrac{1}{2} = \dfrac{3}{16}$
	\end{enumerate}
	\end{multicols}
}


\exe{}{
	Exprimer chaque valeur suivante sous forme de fraction irréductible.
	\begin{multicols}{4}
	\begin{enumerate}[label=\alph*), itemsep=10pt]
		\item $1-\dfrac{4}{44}$
		\item $\dfrac{15}{27}$
		\item $\dfrac{8}{3} + 6$
		\item $\dfrac{65}{5^2}$
		\item $\dfrac{6}{32}\times\dfrac{2}{3}$
		\item $\dfrac{38}{20}-3$
		\item $\left(\dfrac{10}{44}\right)^{-1}$
		\item $\dfrac{6}{18}\times2 + 1$
		\item $\left(\dfrac{66}{99}\right)^2$
		\item $\dfrac{23}{18} + \dfrac12$
		\item $\dfrac{34}{20}$
		\item $2 - \dfrac{9}{5}$
	\end{enumerate}
	\end{multicols}
}{exe:frac-irr}{
	\begin{multicols}{4}
	\begin{enumerate}[label=\alph*), itemsep=10pt]
		\item $\dfrac{10}{11}$
		\item $\dfrac{5}{9}$
		\item $\dfrac{26}{3}$
		\item $\dfrac{13}{5}$
		\item $\dfrac{1}{18}$
		\item $\dfrac{-11}{10}$
		\item $\dfrac{22}{5}$
		\item $\dfrac{5}{3}$
		\item $\dfrac{4}{9}$
		\item $\dfrac{16}{9}$
		\item $\dfrac{17}{10}$
		\item $\dfrac{1}{5}$
	\end{enumerate}
	\end{multicols}
}


\exe{difficulty=0}{
	En décembre pour les fêtes, M. Marchand dit avoir vendu les quatre cinquièmes de sa marchandise. En janvier, pendant les soldes, il a encore vendu les trois quarts de ce qu'il restait.
	\begin{enumerate}
		\item
		Quelle fraction de sa marchandise a-t-il vendu en tout?
		\item  
		La valeur totale de sa marchandise est de 262 000 euros. Quelle somme représente sa vente globale?
	\end{enumerate}

}{exe:Samoyeau6}{
	sol1
}

\exe{difficulty=0}{
	Une usine française de voitures vend $\frac{1}{3}$ de ses voitures en Allemagne, $\frac{1}{4}$ de ce qui reste en Espagne et le reste en France. Quelle proportion de voitures de cette usine reste en France?

}{exe:Samoyeau7}{
	sol1
}

\exe{difficulty=0}{
	Amaya et Mariza ont deux tablettes de chocolat, une au chocolat noir et une au chocolat au lait qui sont de même proportion. 
	Amaya a mangé $\frac{1}{4}$ des $\frac{2}{3}$ de la tablette au chocolat au lait. Mariza a mangé la moitié des $\frac{3}{4}$ de la tablette de chocolat noir. Laquelle a mangé le plus de chocolat? 

}{exe:Samoyeau8}{
	sol1
}

\exe{difficulty=0}{
	Lors d'une course, Pierre fait $\frac{1}{3}$ de la distance totale parcourue en vélo, $\frac{5}{12}$ de la distance totale en courant, $\frac{1}{6}$ de la distance totale à la nage et le reste en kayak. 
	\begin{enumerate}
		\item Compléter la phrase suivante : Pierre a parcouru $\dots\dots$ de la distance totale en kayak.
		\item Cette course fait 24 kilomètres. Quelle est la distance parcourue en kayak? 
	\end{enumerate} 

}{exe:Samoyeau9}{
	sol1
}

\exe{difficulty=0}{
	Considérons le nombre zêta : 
	$ 
	\zeta=\dfrac{ \frac{2}{3} \times \frac{6}{12} + \frac{-3}{7} - 4 }{ \frac{\frac{1}{2}-3}{8\times \frac{-1}{2}} + \frac{2}{5}-\frac{1}{4}}
	$.
	L'exprimer sous forme de fraction irréductible.

}{exe:Samoyeau10}{
	sol1
}

\exe{difficulty=0}{
	Classer les fractions suivantes par ordre croissant : 
	\begin{align*}
		A=\frac{1}{2}, && B=\frac{1}{6}, && C=\frac{5}{12}, &&D=\frac{3}{4}, &&E=\frac{2}{3}.
	\end{align*}
}{exe:Samoyeau2}{
	sol1
}


\exe{}{
	Sans calculatrice, donner le développement décimal des fractions suivantes.
	\begin{multicols}{3}
	\begin{enumerate}[label=\roman*), leftmargin=60pt]
		\item $\dfrac1{10}$
		\item $\dfrac1{5}$
		\item $\dfrac3{10^{5}}$
		\item $\dfrac7{20}$
		\item $\dfrac{395}{50}$
		\item $\dfrac{11}{200}$
	\end{enumerate}
	\end{multicols}
}{exe:dev-decimaux}{
	\begin{multicols}{3}
	\begin{enumerate}[label=\roman*)]
		\item $\dfrac1{10} = 0,1.$
		\item $\dfrac1{5} = \dfrac2{10} = 0,2.$
		\item $\dfrac3{10^{5}} = 0,000~03.$
		\item $\dfrac7{20} = \dfrac{3,5}{10} = 0,35.$
		\item $\dfrac{395}{50} = \dfrac{790}{100} = 7,9.$
		\item $\dfrac{11}{200} = \dfrac{5,5}{100} = 0,055.$
	\end{enumerate}
	\end{multicols}
}


\exe{}{
	Sans calculatrice, donner le développement décimal des fractions suivantes.
	\begin{multicols}{3}
	\begin{enumerate}[label=\roman*), leftmargin=60pt]
		\item $\dfrac1{100}$
		\item $\dfrac3{5}$
		\item $\dfrac2{10^{1}}$
		\item $\dfrac7{2}$
		\item $\dfrac{7}{20}$
		\item $\dfrac{31}{200}$
	\end{enumerate}
	\end{multicols}
}{exe:dev-decimauxE}{
	\begin{multicols}{3}
	\begin{enumerate}[label=\roman*)]
		\item $\dfrac1{100} = 0,01$
		\item $\dfrac3{5} = 0,6$
		\item $\dfrac2{10^{1}} = 0,2$
		\item $\dfrac7{2} = 3,5$
		\item $\dfrac{7}{20} = 0,35$
		\item $\dfrac{31}{200} = 0,155$
	\end{enumerate}
	\end{multicols}
}


\exe{, difficulty=1}{
	Déterminer le développement décimal de $\frac1{11}$ et de $\frac{12}{22}$ en posant la division.
}{exe:111-dec}{
	Posons 
		\[ \dfrac{1}{11} = 0,n_1 n_2 n_3 n_4 \dots. \]
	Par multplication par 10 successives et étude des parties décimales, on obtient $n_1 = 0$, $n_2=9$, et $\frac1{11} = 0,n_3n_4n_5\dots$.
	D'où
		\[ \dfrac1{11} = 0,09~09~09~09~\cdots. \]
	Posons 
		\[ \dfrac{12}{22} = 0,n_1 n_2 n_3 n_4 \dots. \]
	Par multplication par 10 successives et étude des parties décimales, on obtient $n_1 = 5$, $n_2=4$, et $\frac{12}{22} = 0,n_3 n_4 n_5\dots$.
	D'où
		\[ \dfrac{12}{22} = 0,54~54~54~54~\cdots. \]
}



\exe{, difficulty=1}{
	Montrer que $\sqrt2 = 1 + \frac1{1+\sqrt2}$.
	En déduire que $\sqrt2 = 1 + \frac1{2 + \frac1{1+\sqrt2}}$ et que
		\[\sqrt2 = 1 + \cfrac1{2 + \cfrac1{2 + \cfrac1{2 + \cfrac1{1+\sqrt2}}}}. \]
	Exprimer le rationnel $1 + \frac1{2 + \frac1{2 + \frac12}}$ sous la forme $\frac{p}{q}$ avec $p$ et $q$ premiers entre eux et calculer $\left| \sqrt2 - \frac{p}q \right|$ à l'aide de la calculatrice.
	Faire idem avec $1 + \frac1{2 + \frac1{2 + \frac1{2 + \frac12}}}$.
}{exe:frac-continue-sqrt2}{
	En manipulant l'équation on se retrouve à vouloir montrer que $\sqrt2 - 1 = \frac1{\sqrt2+1}$, ce qui est équivalent à $(\sqrt2 - 1)(\sqrt2 + 1) = 1$, vrai par identité remarquable (ou par distributivité simple).
	
	Comme $\sqrt2 = 1 + \frac1{1+\sqrt2}$, on peut remplacer le $\sqrt2$ dans l'expression de droite par toute l'expression de droite.
	On obtient donc $\sqrt2 = 1 + \cfrac1{1 + 1 + \cfrac1{1+\sqrt2}} = 1 + \cfrac1{2 + \cfrac1{1+\sqrt2}}$ comme voulu.
	On obtient l'expression d'après en répétant le même raisonnement.
	
	Par simplifications simples, on trouve $\frac{p}{q} = \frac{17}{12}$.
	La calculatrice nous donne alors $\left| \sqrt2 - \frac{17}{12} \right| \approx 0,002$.
	
	Finalement, remarquons que 
	\begin{align*}
		1 + \cfrac1{2 + \cfrac1{2 + \cfrac1{2 + \cfrac12}}} &= 1 + \cfrac1{1+ 1 + \cfrac1{2 + \cfrac1{2 + \cfrac12}}}, \\
			&= 1 + \cfrac1{1+ \dfrac{17}{12}} = \dfrac{41}{29}.
	\end{align*}
	Le rationnel $\frac{41}{29}$ vérifie $\left| \sqrt2 - \frac{41}{29} \right| \approx 0,0004$, soit 5 fois plus proche de $\sqrt2$ que $\dfrac{17}{12}$ l'est.
}

\exe{difficulty=2}{
	En s'inspirant de l'exercice \ref{exe:frac-continue-sqrt2}, calculer sans calculatrice un rationnel s'approchant de $\sqrt5$ à $10^{-3}$ près.
}{exe:frac-continue-sqrt5}{
	todo
}

\exe{difficulty=1}{
	Vrai ou faux ? Démontrer que chaque proposition est vraie ou donner un contre-exemple.
	\begin{enumerate}
		\item
		La somme de deux nombres rationnels est un nombre rationnel.
		\item
		Le produit de deux nombres rationnels est un nombre rationnel.
		\item
		Un nombre rationnel à la puissance d'un nombre rationnel est un nombre rationnel.
		\item
		Un nombre rationnel à la puissance d'un nombre rationnel est un nombre irrationnel.
		\item
		La somme d'un nombre irrationnel et d'un nombre rationnel est un nombre irrationnel.
		\item
		La somme de deux nombres irrationnels est un nombre irrationnel.
		\item[7. $(\star\star)$]
		Un nombre irrationnel à la puissance d'un nombre irrationnel est irrationnel.
	\end{enumerate}
}{exe:VF-rationnels}{
	todo
}


\exe{, difficulty=1}{
	En musique, la gamme classique contient douze notes avec un rapport de fréquence de 2 sur l'octave.
	La gamme est dite à tempérament égal si le rapport de fréquence entre deux notes consécutives est toujours le même.
	Notons ce rapport $r\in\R$.
	Justifier que $r^{12} = 2$ puis démontrer que $r$ est un nombre irrationnel : $r \not\in\Q$.
}{exe:temp-egal}{
	Le rapport de fréquence entre deux notes consécutives étant toujours $r$, on multiplie $r$ par lui-même 12 fois pour obtenir le rapport de fréquence entre deux notes de distance $12$, c'est-à-dire deux notes séparée par un octave.
	Ainsi $r^{12} = 2$ car l'octave correspond à un rapport de fréquence de 2.
	
	Supposons désormais que $r$ soit rationnel. Alors $r^{6}$ est aussi rationnel.
	Cependant, $(r^6)^2 = r^{12} = 2$, donc $r^{6} = \sqrt2$, qui n'est pas rationnel. \Large\Lightning
}



%%%%%%%%%%%%%%%%%%%%%%%%
%%%%%%%%%% SET 2 %%%%%%%%%%
%%%%%%%%%%%%%%%%%%%%%%%% 



\exe{}{
	Trouver le $x\in\Q$ rationnel vérifiant chaque équation suivante, et l'exprimer sous forme de fraction irréductible.

	\begin{multicols}{2}
	\begin{enumerate}[itemsep=10pt]
		\item $3x + 5 = 2x + 7$
		\item $4x - 6 = 3x + 8$
		\item $5x + 10 = 2x + 4$
		\item $-2x + 3 = 4x - 1$
		\item $7x - 5 = 5x + 9$
		\item $6x + 4 = 3x + 12$
		\item $-3x + 2 = -5x - 4$
		\item $8x - 9 = 6x + 3$
		\item $2x + 11 = -3x + 5$
		\item $-4x + 7 = -x + 10$
	\end{enumerate}
	\end{multicols}
}{exe:inter-affine}
{
	\begin{multicols}{2}
	\begin{enumerate}[itemsep=10pt]
		\item $3x + 5 = 2x + 7 \iff x = 2$
		\item $4x - 6 = 3x + 8 \iff x = 14$
		\item $5x + 10 = 2x + 4 \iff x = -2$
		\item $-2x + 3 = 4x - 1 \iff x = \dfrac{2}{3}$
		\item $7x - 5 = 5x + 9 \iff x = 7$
		\item $6x + 4 = 3x + 12 \iff x = \dfrac{8}{3}$
		\item $-3x + 2 = -5x - 4 \iff x = -3$
		\item $8x - 9 = 6x + 3 \iff x = 6$
		\item $2x + 11 = -3x + 5 \iff x = \dfrac{-6}{5}$
		\item $-4x + 7 = -x + 10 \iff x = -1$
	\end{enumerate}
	\end{multicols}
}



\exe{, difficulty=1}{
	Pour chaque paire d'équations, ajouter entre elles le symbole
		\begin{enumerate}[itemindent=1.5cm]
			\item[$\implies$] si la première équation implique la seconde ;
			\item[$\impliedby$] si la deuxième équation implique la première ;
			\item[$\iff$] si les deux équations sont équivalentes ;
			\item[$\centernot\iff$] si les deux équations sont indépendantes.
		\end{enumerate}
	Spécifier l'opération appliquée pour obtenir l'équation impliquée à partir de l'autre.
	
	
	\begin{center}
	\begin{tabular}{ccc|c|c}
		Équation 1 & Symbole & Équation 2 & Opération & \thead{Les équations sont-\\ elles équivalentes ?} \\ \hline
		$3x + 2y = 1$ &  & $6x  + 4y = 2$ & &  \\ \hline
		$-x + 2y = 1$ &  & $-x + 2y - 1 = 0$ & &  \\ \hline
		$3x + 2y = 1$ & & $-3x - 5y = -3$ & &  \\ \hline
		$x^2 = -3x$ &  & $x = -3$ & &  \\ \hline
		$-10x - y = 3$ &  & $0=0$ & &  \\ \hline
		$x-y= 3$ & & $-x + y = -3$ & &  \\ \hline
		$x -y = 0$ & & $x=y$ & &  \\ \hline
		$-x + 2y = 3$ & & $0=0$ & &  \\ \hline
		$2x + 3y + 4 = 2x + 3y + 5$ & & $1=0$ & &  \\ \hline
		$-x - y = 1$ & & $-x - y = 1$ & &  \\ \hline
		$0=0$ & & $23x - 10y = 6$ & &  \\ \hline
		$x + 2y = 3$ & & $x^2 + 2yx = 3x$ & &  \\ \hline
		$x-3y  =0$ & & $-2x+6y = 0$ & &  \\ \hline
		$-2x + 3y = -3$ & & $4x - 6y = -6$ & &  \\ \hline
		$x^2 = x$ & & $x = 1$ & &  \\ \hline
		$x = 3$ & & $x^2 = 9$ & &  \\ \hline
	\end{tabular}
	\end{center}
	
	{Remarque : par abus de notation et pour gagner du temps, l'absence de symbole signifie en général que les équations sont équivalentes, mais il faut toujours faire attention à ce que cela soit bien le cas !}
	
	\textbf{Attention !} si $f(x) = 0 \implies g(x) = 0$, alors une solution de $f(x) = 0$ est une solution de $g(x) = 0$.
	En revanche, une solution de $g(x) = 0$ peut ne pas être une solution de $f(x) = 0$.
	
}{exe:equations-equivalentes}{
	\begin{center}
	\begin{tabular}{ccc|c|c}
		Équation 1 & Symbole & Équation 2 & Opération & \thead{Les équations sont-\\ elles équivalentes ?} \\ \hline
		$3x + 2y = 1$ & {$\iff$} & $6x  + 4y = 2$ & {$\times2$} & {Oui} \\ \hline
		$-x + 2y = 1$ & {$\iff$} & $-x + 2y - 1 = 0$ & {$-1$} & {Oui} \\ \hline
		$3x + 2y = 1$ & {$\centernot\iff$} & $-3x - 5y = -3$ & & {Non} \\ \hline
		$x^2 = -3x$ & {$\impliedby$} & $x = -3$ & & {Non} \\ \hline
		$-10x - y = 3$ & {$\implies$} & $0=0$ & {$\times0$} & {Non} \\ \hline
		$x-y= 3$ & {$\iff$} & $-x + y = -3$ & {$\times(-1)$} & {Oui} \\ \hline
		$x -y = 0$ &{$\iff$} & $x=y$ & {$+y$} & {Oui} \\ \hline
		$-x + 2y = 3$ & {$\implies$} & $0=0$ & {$\times0$} & {Non} \\ \hline
		$2x + 3y + 4 = 2x + 3y + 5$ & {$\iff$} & $1=0$ & {$+(-2x-3y)$} & {Oui} \\ \hline
		$-x - y = 1$ & {$\centernot\iff$} & $-x - y = 1$ & & {Non} \\ \hline
		$0=0$ & {$\impliedby$} & $23x - 10y = 6$ & & {Non} \\ \hline
		$x + 2y = 3$ & {$\implies$} & $x^2 + 2yx = 3x$ & {$\times x$} & {Non} \\ \hline
		$x-3y  =0$ & {$\iff$} & $-2x+6y = 0$ & {$\times(-2)$} & {Oui} \\ \hline
		$-2x + 3y = -3$ & {$\centernot\iff$} & $4x - 6y = -6$ & & {Non} \\ \hline
		$x^2 = x$ & {$\impliedby$} & $x = 1$ & & {Non} \\ \hline
		$x = 3$ & {$\implies$} & $x^2 = 9$ & {Mise au carré} & {Non} \\ \hline
	\end{tabular}
	\end{center}
}

\exe{difficulty=1}{
	Un élève de Seconde fournit la production écrite suivante.
	Celle-ci présente quelques erreurs.
	
	\fbox{\parbox{\textwidth}{
		\textbf{\oldunderline{Énoncé :}}
		
		Considérons l'équation $x^3 + x^2 + x + 1 = 0$ où $x\in\R$ est un nombre réel encore inconnu.
		\begin{enumerate}
			\item
			Montrer que $(x-1)(x^3 + x^2 + x + 1) = x^4 - 1$.
			\item
			En déduire la ou les solutions de l'équation originelle.
		\end{enumerate}
		
		\textbf{\underline{Réponse de l'élève :}}
		
		\begin{enumerate}
			\item
			Pour développer $(x-1)(x^3 + x^2 + x + 1)$,
			je distribue d'abord le $x$, puis le $-1$, et j'ajoute les résultats.
				\[ x(x^3 + x^2 + x + 1) + (-1)(x^3 + x^2 + x + 1) = x^4 + x^3 + x^2 + x -x^3 -x^2 - x - 1. \]
			Les puissances de $x$ égales se regroupent pour simplifier l'expression
				\[ x^4 + x^3 + x^2 + x -x^3 -x^2 - x - 1 = x^4 - 1, \]
			comme voulu.
			\item
			Un $x\in\R$ pour lequel $x^3 + x^2 + x + 1 = 0$ vérifie donc $0 = x^4 - 1$ et donc
			$x^4  = 1$.
			En appliquant la racine carrée, je trouve $x^2 = 1$ car $x^4 = (x^2)^2$.
			En l'appliquant encore, j'obtiens finalement $x = 1$.
			
			L'unique solution de l'équation est donc $1$.
		\end{enumerate}
	}}
	\begin{enumerate}
		\item 
		Identifier les erreurs commises par l'élève.
		\item
		Répondre à l'énoncé.
	\end{enumerate}
}{exe:erreurs0}{
	\begin{enumerate}
		\item
		La première question est bien traitée et ne comporte aucune erreur.	
		\item
		La deuxième question contient deux erreurs : une lors de la résolution de $x^4 = 1$, et l'autre dans sa conclusion.
		
		Premièrement, appliquer la racine carrée est pertinente, mais $\sqrt{x^2} = |x|$, et non pas $x$ !
		La valeur absolue est importante pour ne pas oublier de solution.
		Ainsi, si $x^4 = 1$, alors $|x^2| = 1$. Comme $x^2 \geq 0$ pour tout $x\in\R$, la valeur absolue peut disparaître ici.
		Puis, si $x^2 = 1$, alors $|x| = 1$, et donc $x=1$ ou $x=-1$. Il y a donc deux solutions réelles à l'équation $x^4  = 1$.
		
		Deuxièmement, la conclusion est erronnée car elle présente une erreur de logique.
		En effet, l'impliquation $x^3 + x^2 + x + 1 = 0 \implies x \in \{-1 ; 1\}$ est vraie, mais l'implication inverse ne l'est pas !
		Nous disposons de deux candidats de solutions qu'il faut vérifier. La solution $x=1$ ne fonctionne pas, mais $x=-1$, elle, fonctionne.
		La seule solution est donc $x=-1$.
		
		En multipliant l'équation originelle par $x-1$, nous avons ajouté artificiellement la solution $x=1$.
		Par exemple, les implications suivantes sont vraie
			\[ x - 1 = 0 \implies (x-3)(x-1) = 0 \implies x \in \{3 ; 1\}. \]
	\end{enumerate}
}

\exe{}{
	Résoudre les équations suivantes pour $x\in\R$ non nul. 
	%Exprimer $x$ sous la forme $q \sqrt{n}$ avec $q\in\Q$ rationnel et $n\in\N$ le plus petit possible.
	\begin{multicols}{4}
	\begin{enumerate}
		\item $\dfrac{x}3 = \dfrac12.$
		\item $ \dfrac7x  =\dfrac12$
		\item $\dfrac6x  =\dfrac{\sqrt{3}}2$
		\item $\dfrac{x}2 =\dfrac1{\sqrt{3}}$
		%\item $\dfrac9x  =\dfrac{\sqrt{6}}5$
		%\item $-\dfrac3x =\sqrt{7}$
		%\item $-\dfrac7{2x}  =\dfrac{\sqrt{11}}3$
		%\item $\dfrac3{5x} =-\sqrt{5}$
	\end{enumerate}
	\end{multicols}
}{exe:eq-inv-x}{
	\begin{multicols}{2}
	\begin{enumerate}
		\item 
			\begin{align*}
				\dfrac{x}3 &= \dfrac12 \\
				x &= 3 \cdot \frac12 \\
				x &= \frac32
			\end{align*}
		\item
			\begin{align*}
				\dfrac7x  &=\dfrac12 \\
				7 &= x \cdot \frac12 \\
				x &= 7\cdot2 = 14
			\end{align*}
		\item
			\begin{align*}
				\dfrac6x  &= \dfrac{\sqrt{3}}2 \\
				6 &= x \cdot \dfrac{\sqrt{3}}2 \\
				x &= 6\cdot\dfrac2{\sqrt{3}} \\
				x &= \frac{12}3 \sqrt3 = 4\sqrt3
			\end{align*}
		\item
			\begin{align*}
				\dfrac{x}2 &= \dfrac1{\sqrt{3}} \\
				x &= 2 \cdot \dfrac1{\sqrt{3}} \\
				x &= \frac23 \sqrt3
			\end{align*}
	\end{enumerate}
	\end{multicols}
}

\exe{}{
	Trouver le $x \in \Q$ rationnel (et tel que chaque dénominateur est non nul) vérifiant chaque équation suivante, et l'exprimer sous forme de fraction irréductible.
	
	\begin{multicols}{3}
	\begin{enumerate}[itemsep=10pt]
		\item $9 = \dfrac3x$
		\item $-4 = \dfrac2x$
		\item $0,2422 = \dfrac1x$
		\item $0,7 = \dfrac{0,1}x$
		\item $3 - \dfrac4x = -9 + \dfrac1x$
		\item $-2 + \dfrac{-2}x = \dfrac3x - 10$
		\item $\dfrac1{x-2} = 3$
		\item $\dfrac3{x+5} = \dfrac7{x+5} - 3$
	\end{enumerate}
	\end{multicols}

}{exe:eq-lin0}
{
	\begin{multicols}{2}
	\begin{enumerate}[itemsep=10pt]
		\item $9 = \dfrac3x \iff x = \dfrac13$
		\item $-4 = \dfrac2x \iff x=-\dfrac12$
		\item $0,2422 = \dfrac1x \iff x = \dfrac{10000}{2422} = \dfrac{5000}{1211}$
		\item $0,7 = \dfrac{0,1}x \iff x = \dfrac17$
		\item $3 - \dfrac4x = -9 + \dfrac1x \iff x = \dfrac5{12}$
		\item $\dfrac1{x-2} = 3 \iff x=\dfrac73$
		\item $\dfrac3{x+5} = \dfrac7{x+5} - 3 \iff x=\dfrac{-11}3$
	\end{enumerate}
	\end{multicols}
}


\exe{difficulty=1}{
	Montrer que les nombres suivants sont rationnels en les exprimant sous forme de fraction d'entiers.
	\[
	\begin{aligned}
		A &= 0,666{\dots} \\
		B &= 1,666{\dots} \\
	\end{aligned}
	\hspace{5cm}
	\begin{aligned}
		C &= 0,121212{\dots} \\
		D &= 0,34777{\dots} \\
	\end{aligned}
	\]
	\[
	E = 0,123456789123456789123{\dots} \text{ (nombre d'Amandine) }
	\]
	
	\emph{Indication : montrer d'abord que $10A = 6+A$.}
}{exe:dev-to-fraction}{
	\begin{enumerate}[label=\Alph*.]
		\item 
		La relation $10A = 6 + A$ donne $A = \frac23$.
		\item
		La relation $10B = 15 + B$ donne $B = \frac{15}{9} = \frac53$.
		On aurait aussi pû utiliser que $B = 1+A$.
		\item
		La relation $100C = 12 + C$ donne $C = \frac{12}{99} = \frac{4}{33}$.
		\item
		En posant $\tilde{E} = 100E = 34,777{\dots}$, on obtient $10\tilde{E} = 313 + \tilde{E}$, d'où $\tilde{E} = \frac{313}{9}$, et donc $E = \frac{1}{100}E' = \frac{313}{900}$.
		\item 
		La relation $10^9 E = 123~456~789 + E$ donne $E = \frac{123~456~789}{999~999~999}$.
	\end{enumerate}
}

\exe{}{
	Montrer que les nombres suivants sont rationnels en les exprimant sous forme de fraction d'entiers.
	\begin{align*}
		A = 0,123123123{\dots} &&
		B = 0,9999{\dots}
	\end{align*}
}{exe:dev-to-fractionE}{
	\begin{enumerate}[label=\Alph*.]
		\item
		La relation $1000A = 123 + A$ donne $A = \frac{123}{999} = \frac{41}{333}$.
		\item 
		La relation $10B = 9 + B$ donne $9B = 9$ et $B=1$.
	\end{enumerate}
}


\exe{, difficulty=1}{
	Multipliez l'équation 
		\[ \dfrac{a}b + \dfrac{c}d = x \]
	par $b$ puis $d$ à gauche et à droite et en déduire la règle d'addition des fractions :
		\[ x = \dfrac{ad + cb}{bd}. \]
}{exe:common-denom}{
	En multipliant par $bd$, on trouve
		\[ (bd) x =  bd \left( \dfrac{a}b + \dfrac{c}d \right) = ad + bc, \]
	d'où
		\[ x = \dfrac{ad+bc}{bd}. \]
}

\exe{}{
	Résoudre les équations suivantes pour $x\in\R$ non nul. Exprimer $x$ sous la forme $q \sqrt{r}$ avec $q\in\Q$ rationnel et $r\in\N$ le plus petit possible.
	\begin{multicols}{2}
	\begin{enumerate}
		\item $\dfrac3x = \dfrac73.$
		\item $ \dfrac7x  =\dfrac1x + 2$
		\item $\dfrac{16}x  =\dfrac{\sqrt{3}}2$
		\item $-\dfrac2x =4\sqrt{3}$
		\item $\dfrac9{-x}  =\dfrac{\sqrt{6}}5$
		\item $-\dfrac3x =\dfrac57\sqrt{7}$
		\item $-\dfrac9{2x}  =\dfrac{\sqrt{11}}3$
		\item $\dfrac2{7x} =-\sqrt{14}$
	\end{enumerate}
	\end{multicols}
}{exe:eq-sqrt}{}



\exe{}{
	On souhaite connaître les solutions de l'équation du troisième degré d'inconnue $x\in\R$ :
		\begin{align}
			4x^3 - 6x^2 - 2x + 3  = 0. \label{eq:2}
		\end{align}
	\begin{enumerate}
		\item Montrer que $4x^3 - 6x^2 - 2x + 3 = (2x-3) \left(2x^2-1\right)$.
		\item En déduire l'ensemble des solutions de l'équation \eqref{eq:2}.
	\end{enumerate}
}{exe:eq8}{

	\begin{enumerate}
		\item On part toujours de la forme factorisée qu'on développe par double distributivité.
		On utilise aussi que $x^2 \cdot x = x \cdot x \cdot x = x^3$.
			\[ (2x - 3)(2x^2 - 1) = 4x^3 - 2x -6x^2 + 3. \]
		\item On utilise que le produit est nul si et seulement si un des deux facteurs est nul.
		Or l'équation $2x^2 - 1 = 0$ est équivalente à $x^2 = \dfrac12$.
		En prenant la racine et en utilisant que $\sqrt{x^2} = |x|$, on obtient
			\[ |x| = \sqrt{\dfrac12} = \dfrac1{\sqrt2}. \]
		Les deux valeurs $\pm \dfrac1{\sqrt2}$ (notation $\pm$ pour \og plus ou moins \fg) sont donc solutions.
		On en déduit que l'ensemble des solutions de \eqref{eq:2} est $\left\{ \dfrac32 ; \pm \dfrac1{\sqrt2} \right\} = \left\{ \dfrac32 ; \dfrac1{\sqrt2} ; - \dfrac1{\sqrt2} \right\}$
	\end{enumerate}

}


\exe{}{
	On souhaite connaître les solutions de l'équation du quatrième degré d'inconnue $x\in\R$ :
		\begin{align}
			16x^4 - 24x^2 + 9  = 0. \label{eq:3}
		\end{align}
	\begin{enumerate}
		\item Montrer que $16x^4 - 24x^2 + 9 = \left(4x^2-3\right)^2$.
		\item En déduire l'ensemble des solutions de l'équation \eqref{eq:3}.
	\end{enumerate}
}{exe:eq9}{

	\begin{enumerate}
		\item On part toujours de la forme factorisée qu'on développe à l'aide de l'identité remarquable $(a-b)^2 = a^2 + b^2 -2ab$.
		On utilise aussi que $(x^2)^2 = x^2 \cdot x^2 = x\cdot x\cdot x\cdot x = x^4$.
			\[ (4x^2 - 3)^2 = (4x^2)^2 + 3^2 - 2\cdot(4x^2)\cdot3 = 16x^4 + 9 -24x^2. \]
		\item On utilise que le produit est nul si et seulement si un des deux facteurs est nul.
		Ici, les deux facteurs sont les mêmes, car $(4x^2 - 3)^2 = (4x^2 - 3)(4x^2 - 3)$, donc on se remet à résoudre $(4x^2 - 3) = 0$.
		Ceci est équivalent à $x^2 = \dfrac34$, et donc l'ensemble des $x$ vérifiant \eqref{eq:3} est donné par $\left\{ \pm \sqrt{\dfrac34} \right\} = \left\{ \pm \dfrac12\sqrt{3} \right\}= \left\{ \dfrac12\sqrt{3} ;  -\dfrac12\sqrt{3} \right\}$.
	\end{enumerate}

}




%%%%%%%%%%%%%%%%%%%%%%%%
%%%%%%%%%% SET 2 %%%%%%%%%%
%%%%%%%%%%%%%%%%%%%%%%%% 


\exe{}{
	On souhaite connaître les solutions de l'équation du deuxième degré d'inconnue $x\in\R$ :
		\begin{align}
			22x^2 - 125x + 22  = 0. \label{eq:1}
		\end{align}
	\begin{enumerate}
		\item Montrer que $22x^2 - 125x + 22 = (2x-11)(11x-2)$.
		\item En déduire l'ensemble des solutions de l'équation \eqref{eq:1}.
	\end{enumerate}
}{exe:eq7}{
	\begin{enumerate}
		\item On part toujours de la forme factorisée qu'on développe par double distributivité :
			\[ (2x - 11)(11x - 2) = 22x^2 - 4x - 121x + 22 = 22x^2 - 125x + 22. \]
		\item On utilise que le produit est nul si et seulement si un des deux facteurs est nul.
		L'ensemble des solutions de \eqref{eq:1} est donc $\left\{ \dfrac{11}2 ; \dfrac2{11} \right\}$.
	\end{enumerate}
}


\exe{}{
	Pour chaque propriété, donner une fonction polynomiale $f$ à coefficients réels et non nulle la vérifiant.
	\begin{multicols}{2}
	\begin{enumerate}
		\item $f$ s'annule en $1$.
		\item $f$ s'annule en $-10$.
		\item $f$ s'annule en $0$.
		\item $f$ s'annule en $-10$ et en $1$.
		\item Les racines de $f$ sont $2, -3, \frac27$, et $0$.
	\end{enumerate}
	\end{multicols}
}{exe:f-zeroes}{
	Le polynôme identiquement nul $f(x) = 0$ n'est bien sûr pas autorisé sinon l'exercice n'a pas d'intérêt !
	\begin{enumerate}
		\item $f(x) = x-1$ fonctionne ainsi que $f(x) = 3(x-1)$, ou $f(x) = -(x-1) = 1-x$.
		\item $f(x) = x+10$ fonctionne, ainsi que $f(x) = (x+10)(x-1)$, et tous ses multiples.
		\item $f(x) = x$ ou $f(x) = 500x$.
		\item $f(x) = (x+10)(x-1) = x^2 +9x -10$.
		\item $f(x) = (x-2)(x+3)(x-\frac27)x$.
	\end{enumerate}

}

\exe{}{
	Résoudre les équations suivantes pour $x\in\R$ non nul. Exprimer $x$ sous la forme $q \sqrt{r}$ avec $q\in\Q$ rationnel et $r\in\N$ le plus petit possible.
	\begin{multicols}{2}
	\begin{enumerate}
		\item $\dfrac3x = \dfrac73.$
		\item $ \dfrac7x  =\dfrac1x + 2$
		\item $\dfrac{16}x  =\dfrac{\sqrt{3}}2$
		\item $-\dfrac2x =4\sqrt{3}$
		\item $\dfrac9{-x}  =\dfrac{\sqrt{6}}5$
		\item $-\dfrac3x =\dfrac57\sqrt{7}$
		\item $-\dfrac9{2x}  =\dfrac{\sqrt{11}}3$
		\item $\dfrac2{7x} =-\sqrt{14}$
	\end{enumerate}
	\end{multicols}
}{exe:eq-sqrt}{}




\exe{}{
	On souhaite connaître les solutions de l'équation du troisième degré d'inconnue $x\in\R$ :
		\begin{align}
			4x^3 - 6x^2 - 2x + 3  = 0. \label{eq:2}
		\end{align}
	\begin{enumerate}
		\item Montrer que $4x^3 - 6x^2 - 2x + 3 = (2x-3) \left(2x^2-1\right)$.
		\item En déduire l'ensemble des solutions de l'équation \eqref{eq:2}.
	\end{enumerate}
}{exe:eq8}{

	\begin{enumerate}
		\item On part toujours de la forme factorisée qu'on développe par double distributivité.
		On utilise aussi que $x^2 \cdot x = x \cdot x \cdot x = x^3$.
			\[ (2x - 3)(2x^2 - 1) = 4x^3 - 2x -6x^2 + 3. \]
		\item On utilise que le produit est nul si et seulement si un des deux facteurs est nul.
		Or l'équation $2x^2 - 1 = 0$ est équivalente à $x^2 = \dfrac12$.
		En prenant la racine et en utilisant que $\sqrt{x^2} = |x|$, on obtient
			\[ |x| = \sqrt{\dfrac12} = \dfrac1{\sqrt2}. \]
		Les deux valeurs $\pm \dfrac1{\sqrt2}$ (notation $\pm$ pour \og plus ou moins \fg) sont donc solutions.
		On en déduit que l'ensemble des solutions de \eqref{eq:2} est $\left\{ \dfrac32 ; \pm \dfrac1{\sqrt2} \right\} = \left\{ \dfrac32 ; \dfrac1{\sqrt2} ; - \dfrac1{\sqrt2} \right\}$
	\end{enumerate}

}


\exe{}{
	On souhaite connaître les solutions de l'équation du quatrième degré d'inconnue $x\in\R$ :
		\begin{align}
			16x^4 - 24x^2 + 9  = 0. \label{eq:3}
		\end{align}
	\begin{enumerate}
		\item Montrer que $16x^4 - 24x^2 + 9 = \left(4x^2-3\right)^2$.
		\item En déduire l'ensemble des solutions de l'équation \eqref{eq:3}.
	\end{enumerate}
}{exe:eq9}{

	\begin{enumerate}
		\item On part toujours de la forme factorisée qu'on développe à l'aide de l'identité remarquable $(a-b)^2 = a^2 + b^2 -2ab$.
		On utilise aussi que $(x^2)^2 = x^2 \cdot x^2 = x\cdot x\cdot x\cdot x = x^4$.
			\[ (4x^2 - 3)^2 = (4x^2)^2 + 3^2 - 2\cdot(4x^2)\cdot3 = 16x^4 + 9 -24x^2. \]
		\item On utilise que le produit est nul si et seulement si un des deux facteurs est nul.
		Ici, les deux facteurs sont les mêmes, car $(4x^2 - 3)^2 = (4x^2 - 3)(4x^2 - 3)$, donc on se remet à résoudre $(4x^2 - 3) = 0$.
		Ceci est équivalent à $x^2 = \dfrac34$, et donc l'ensemble des $x$ vérifiant \eqref{eq:3} est donné par $\left\{ \pm \sqrt{\dfrac34} \right\} = \left\{ \pm \dfrac12\sqrt{3} \right\}= \left\{ \dfrac12\sqrt{3} ;  -\dfrac12\sqrt{3} \right\}$.
	\end{enumerate}

}


\exe{difficulty=1}{
	\begin{multicols}{2}
	La Terre complète son tour autour du Soleil en environ 365,2422 jours. 
	Un calendrier ne peut donc pas contenir 365 jours tous les ans si celui-ci doit être synchronisé avec les saisons.
	
	Si le calendrier contient moins de 365,2422 jours, il sera en avance.
	S'il contient plus de 365,2422 jours, il sera en retard

	\begin{center}
		\includegraphics[page=4]{../../CM/2-figures.pdf}
	\end{center}
	\end{multicols}
	
	\begin{enumerate}
	\item
	Calendrier julien\footnote{De Jules César (né en 100 av. J.-C., mort en 44 av. J.-C.).}.
	
	En $46$ av. J.-C., Jules César introduit la règle de l'année bissextile.
	\begin{enumerate}
		\item
		Trouver un entier naturel $a \in \N$ non nul tel que
			\[ 365{,}2422 \approx 365 + \dfrac1a.\]

		\item
		En déduire la règle d'année bissextile du calendrier julien.
		
		\item
		Au bout de 1 600 ans, de combien de jours le calendrier est-il en retard ?
	\end{enumerate}
	
	\item
	Calendrier grégorien\footnote{Du pape Grégoire XIII (1502-1585).}.
	
	En 1582, plus de 1 600 ans après l'introduction de la règle d'année bissextile, le calendrier julien prend donc un retard notable.
	Le pape Grégoire XIII promulgue alors une réforme du calendrier, ainsi créant le calendrier grégorien encore utilisé aujourd'hui.
	
	Afin de corriger le problème de dérive du calendrier, une nouvelle règle est ajoutée à la définition d'année bissextile.
	
	\begin{enumerate}
		\item
		Trouver un entier naturel $b \in \N$ tel que
			\[ 365{,}2422 \approx 365 + \dfrac14 - \dfrac{b}{400}.\]

		\item
		En déduire la règle d'année bissextile du calendrier grégorien.
		
		\item
		Au bout de combien d'années le calendrier est-il en retard d'un jour ?
	\end{enumerate}
	
	\item
	Un nouveau calendrier ?
	\begin{enumerate}
		\item
		Trouver un entier naturel $c \in \N$ tel que
			\[ 365{,}2422 \approx 365 + \dfrac14 - \dfrac{3}{400} - \dfrac{c}{4000}.\]

		\item 
		Créer une nouvelle règle d'année bissextile.
		\item
		Calculer le nombre d'années avant que votre calendrier soit en retard d'un jour.
	\end{enumerate}
	\end{enumerate}
}{exe:calendrier}{}





\exe{, difficulty=2}{
	Donner quelques termes en plus de la somme infinie suivante. 
		\[ \dfrac1{10} + \dfrac1{100} + \dfrac1{1 000} + \dfrac1{10 000} + \dots \]
	Écrire le nombre en écriture décimale, puis l'exprimer sous forme de fraction de deux entiers comme à l'exercice \ref{exe:dev-to-fraction}.
}{exe:20}{
	La somme infinie vaut $0,1111{\dots} = \frac19$.
}

\exe{, difficulty=2}{
	Donner quelques termes en plus de la somme infinie suivante.
		\[ \dfrac12 + \dfrac14 + \dfrac18 + \dfrac1{16} + \dots \]
	À quoi est-elle égale ? Un résultat entier est attendu.
}{exe:21}{
	Posons $S$ cette somme. 
	En multipliant par $2$, on trouve $2S = 1 + S$, et donc $S=1$.
}

\exe{, difficulty=2}{
	On considère la somme finie suivante
		\[ S = 1 + 2 + 4 + 8 + 16 + \dots + 2^{n-1} +  2^n, \]
	où $n\in\N$ est un entier naturel.
	
	Montrer que $2S = S - 1 + 2^{n+1}$, et en déduire la valeur de $S$ sous forme close (sans points de suspension).
}{exe:22}{
	En multipliant par 2, le premier terme (1) disparait, et un nouveau terme ($2^{n+1}$) apparaît.
	On a donc bien $2S = S -1 + 2^{n+1}$, et donc $S = 2^{n-1} - 1$.
}


\exe{, difficulty=1}{
	Montrer qu'un triangle isocèle rectangle ne peut pas avoir des côtés de longueurs toutes entières.
}{exe:non-iso-rec}{
	Supposons qu'il existe un triangle isocèle rectangle de côtés $(a ; a; b)$.
	
	Or, comme le triangle est rectangle isocèle,
		\[ b^2 = a^2 + a^2 = 2a^2 \iff b = \pm\sqrt2 a. \]
	Par conséquent, $\sqrt2 = \pm\frac{b}{a} \in \Q$, ce qui contredit le théroème d'irrationalité de la racine carrée de 2.
}

\exe{, difficulty=1}{
	Soit $q\in\Q$ un rationnel. Montrer qu'un triangle dont les côtés sont de longueur $(4 ; 5 ; q)$ ne peut être rectangle que si $q=3$.
	
	En autorisant $x\in\R$ irrationel, donner un triangle rectangle de côtés $(4 ; 5; x)$ avec $x\neq3$.
}{exe:non-rec}{
	Il y a deux cas de figure à étudier : $q$ est la longueur de l'hypothénuse, ou elle ne l'est pas.
	Dans le premier cas, le théorème de Pythagore énonce que
		\[ 4^2 + 5^2 = q^2 \iff 41 = q^2. \]
	Il suit que $q=\sqrt{41}$ qui n'est pas rationnel car 41 n'est pas un carré parfait.
	Ceci répond à la deuxième question : $(4 ; 5; \sqrt{41})$ est rectangle.
	
	Dans le second cas,
		\[ q^2 + 4^2 = 5^2 \iff q^2 = 9. \]
	Par conséquent, $q = \pm3$ et seule la solution $q=3$ est positive et donc une longueur valide.
}
